Let $X_{2} = T_{3}$ (the third generator of $\su(2)$, eigenvalues $\pm\tfrac12$ on the doublet), $X_{3}$ the analogous generator of $\su(3)$ (eigenvalues $\pm\tfrac12$ on two colours, $0$ on the third) and $X_{1} = \sqrt{3/5}\,Y$, three elements of $\su(5)$ through the embedding in which the fundamental $5$ restricts to $(2,1)_{1/2} \oplus (1,3)_{-1/3}$. On the $5$: $\tr X_{2}^{2} = 2 \cdot \tfrac14 = \tfrac12$; $\tr X_{3}^{2} = 2 \cdot \tfrac14 = \tfrac12$; $\tr X_{1}^{2} = \tfrac35\,(2 \cdot \tfrac14 + 3 \cdot \tfrac19) = \tfrac35 \cdot \tfrac56 = \tfrac12$. By Fact \claimref{standard-model.imported-invariant-forms-simple-lie-algebra} the form $\tr(\rho(X)\rho(Y))$ on $\su(5)$ is $c_{\rho}$ times the Killing form for every representation $\rho$; since the three elements have equal norm $\tfrac12$ in the form of the $5$, they have equal Killing norm, hence equal norm $c_{\rho}\,\kappa$ in the form of every $\rho$. The index $T_{i}(\rho)$ is by definition $\tr_{\rho}(X_{i}^{2})$ in the normalisation of Fact \claimref{standard-model.imported-coefficient-sum-indices} (the fundamental's value $\tfrac12$ for the non-abelian factors, and $\tfrac35 \sum Y^{2}$ for the abelian one, which is $\tr_{\rho}X_{1}^{2}$), so $T_{1}(\rho) = T_{2}(\rho) = T_{3}(\rho)$. A representation of $\mathrm{Spin}(10)$ restricts to a representation of $SU(5)$ and the same conclusion holds; a direct sum of representations of $G_{U}$ is one. The shift then follows from Fact \claimref{standard-model.imported-coefficient-sum-indices} with the common threshold making the content active over one interval. The script checks the equality on the Weyl $10$ ($\Dl = (1,1,1)$), the scalar $\bar{5}$ ($\tfrac16\,(1,1,1)$), the Weyl $24$ ($\tfrac{10}{3}\,(1,1,1)$, against the gauginos alone at $(0, \tfrac43, 2)$), one Weyl generation ($\tfrac43\,(1,1,1)$) and three scalar generations ($2\,(1,1,1)$).
